\section{Encoded Paper Catalogue}
\label{sec:appendix-paper-catalogue}
%% ====================================================================

This appendix catalogues a corpus of recently published MLIP studies
encoded in the \MLIPsOntology{}. The first subsection
(\S\ref{sec:appendix-example-protocol}) defines the extraction
protocol used to produce the per-paper reports; the remaining
subsections apply that protocol to one paper each. The corpus
underpins the competency-question evaluation in
\S\ref{sec:eval-cq}: the seeded knowledge graph used by the SPARQL
templates is the union of the per-paper canonical Turtle files
listed below.

The detailed Kumar~et~al.\ (2025) worked example in
\S\ref{sec:appendix-example-kumar2025} is the prose-rich reference
for what each encoding aspires to; the catalogue entries below
follow the leaner, protocol-driven format.

\subsection{Extraction Protocol}
\label{sec:appendix-example-protocol}

The Kumar~et~al.\ worked example
(\S\ref{sec:appendix-example-kumar2025}) is the reference for what
an encoded MLIP paper should look like. To replicate that quality
across the catalogue at scale---and to give a future agentic-AI
tool something concrete to follow---we fix a 12-question extraction
protocol. The protocol has two artefacts per paper: a canonical
Turtle file carrying the encoded data, and a subsection in this
appendix carrying the prose. The two are kept in sync by an
automated \emph{round-trip check} described at the end of this
subsection.

\paragraph{Overview of the flow.}
For each paper we ask the same 12 questions in fixed order. Each
question (a) produces a Turtle fragment that is appended to the
paper's canonical file
\texttt{artifacts/kg/papers/\textit{paper-id}.ttl}, and (b) a short
commentary paragraph in the paper's subsection. Each of questions
Q1--Q11 has an associated CONSTRUCT query that, when run against the
paper's canonical file, regenerates the Turtle listing shown under
that question's commentary. The union of the 11 CONSTRUCT outputs is
required (modulo blank-node renaming and prefix declarations) to
reproduce the canonical file---this is the round-trip check. Q12 is
prose-only (gaps are not always reducible to triples) and is excluded
from the round-trip.

\paragraph{The 12 questions.}
Grouped into five phases:

\smallskip
\noindent\textbf{Phase A.\ Identification.}
\begin{itemize}
\item[Q1.] What are the bibliographic details of the source paper?
  Title, authors, year, venue, DOI.
  \\ \emph{Targets:} $\dlConcept{schema:ScholarlyArticle}$,
  $\BenchmarkStudyC$, $\reportedInR$.
\end{itemize}

\smallskip
\noindent\textbf{Phase B.\ Training data.}
\begin{itemize}
\item[Q2.] What material system(s) are studied? Chemical formula,
  material class (alloy, intermetallic, Laves phase, molecular,
  oxide), structural specifics (phase, space group, defect type),
  Wikidata link if available.
  \\ \emph{Targets:} $\MaterialSystemC$ and its data properties.
\item[Q3.] What method produced the reference data? DFT,
  wave-function (CCSD(T), MP2, CASPT2, \dots), AIMD, or experimental.
  \\ \emph{Targets:} $\ReferenceCalculationC$ and its concrete
  subtype.
\item[Q4.] What reference-method settings are reported? For DFT:
  $\xcFunctionalR$, $\kPointMeshR$, $\energyCutoffR$,
  $\pseudopotentialTypeR$. For wave-function: $\wfMethodR$,
  $\basisSetR$, $\frozenCoreR$. For both: software code via
  $\dlRole{usedReferenceCode}$ (or its DFT-specific sub-property).
  \\ \emph{Targets:} $\ReferenceSettingsC$ subtypes and $\LibraryC$.
\item[Q5.] How is the training dataset characterised? Number of
  configurations, covered physical properties (energies, forces,
  stresses, virials), provenance class, link to the reference
  calculation.
  \\ \emph{Targets:} $\TrainingDatasetC$, $\DatasetProvenanceC$,
  $\CoveredPropertyC$, $\hasReferenceCalculationR$.
\item[Q6.] What sampling strategies were used? Chemical, vibrational
  or MD-driven, active-learning, perturbation-based, exhaustive
  enumeration, or combinations.
  \\ \emph{Targets:} $\SamplingStrategyC$ instances and
  $\samplingStrategyR$.
\end{itemize}

\smallskip
\noindent\textbf{Phase C.\ Method, run, and model.}
\begin{itemize}
\item[Q7.] What MLIP method is used, and what are its components?
  Functional form, loss function, training algorithm (the ML-Schema
  alignment point), software implementation and version.
  \\ \emph{Targets:} $\dlConcept{MLIPMethod}$,
  $\dlConcept{FunctionalForm}$, $\dlConcept{LossFunction}$,
  $\hasTrainingAlgorithmR$, $\ImplementationC$, $\LibraryC$.
\item[Q8.] What hyperparameter settings does the paper explicitly
  report? Cutoff radius, MTP level, number of layers, learning rate,
  etc. Implicit defaults are recorded as
  \emph{not reported} but referenced in Q12.
  \\ \emph{Targets:} $\HyperparameterSettingC$ and
  $\forHyperparameterR$.
\item[Q9.] What is the structure of the run? Single training run vs.\
  multi-step active-learning pipeline; one trained model vs.\ a per-
  phase or per-fold ensemble.
  \\ \emph{Targets:} $\dlConcept{MLIPRun}$, $\TrainedModelC$,
  $\appliesMethodR$, $\runsOnR$, $\producesR$,
  $\hasHyperparameterSettingR$.
\end{itemize}

\smallskip
\noindent\textbf{Phase D.\ Evaluation.}
\begin{itemize}
\item[Q10.] What benchmark results does the paper report? Each
  accuracy metric is recorded with its type (RMSE, MAE, \dots), the
  property it measures (energy, forces, stresses, downstream
  properties), and its numerical value with units.
  \\ \emph{Targets:} $\BenchmarkResultC$, $\AccuracyMetricC$,
  $\MetricTypeC$, $\MetricPropertyC$, $\evaluatesModelR$,
  $\targetMaterialR$.
\item[Q11.] What computational-resource metadata is reported?
  Training duration, GPU hours, training hardware, peak memory;
  inference time per atom, inference hardware.
  \\ \emph{Targets:} data properties on $\dlConcept{MLIPRun}$ and
  $\TrainedModelC$.
\end{itemize}

\smallskip
\noindent\textbf{Phase E.\ Gaps.}
\begin{itemize}
\item[Q12.] What concepts that the ontology can express does the
  paper \emph{not} report? Free-form, prose-only. This question
  produces no Turtle and is excluded from the round-trip check.
\end{itemize}

\paragraph{Per-paper subsection template.}
A protocol-derived subsection contains 12 paragraphs corresponding
to Q1--Q12 (in order). Each paragraph for Q1--Q11 has the structure:
(i) a 2--4 sentence prose answer derived from the source paper;
(ii) a Turtle listing produced by the question's CONSTRUCT query;
and (iii) where the paper does not report a piece of data the
ontology can model, an explicit \emph{not reported} marker so that
Q12 has a complete inventory to summarise. Q12 is prose-only.

\paragraph{Files and naming.}
Per-paper data and queries live under \texttt{artifacts/kg/}:

\begin{itemize}
\item \texttt{papers/\textit{paper-id}.ttl}---canonical Turtle for
  one paper. Identifiers use a \texttt{firstauthor-year} convention
  (\texttt{kumar2025.ttl}, \texttt{qi2023.ttl}, \dots).
\item \texttt{queries/q01-bibliographic.rq} \dots
  \texttt{queries/q11-resources.rq}---one CONSTRUCT query per
  question Q1--Q11. The queries are paper-agnostic; they are
  parameterised only by the prefix declarations of the paper they
  are run against.
\item \texttt{check-roundtrip.sh \textit{paper-id}}---runs the 11
  CONSTRUCT queries against
  \texttt{papers/\textit{paper-id}.ttl}, takes their union, and
  diff-canonicalises against the canonical file. Exits non-zero on
  drift.
\item \texttt{build-listings.sh \textit{paper-id}}---runs the 11
  queries and emits Turtle listings ready for inclusion in the
  paper's subsection.
\end{itemize}

\paragraph{Round-trip check.}
At the end of every paper's encoding pass, we run
\texttt{check-roundtrip.sh}. The check parses the 11 query outputs
into N-triples (via \texttt{rapper -o ntriples}), sorts and dedups,
and diffs against the same canonicalisation of the paper's
\texttt{.ttl} file. Two outcomes are possible:

\begin{itemize}
\item \emph{Pass.} The 11 questions partition the paper's data; the
  subsection is a faithful presentation of the canonical file.
\item \emph{Fail.} A triple is in the canonical file but not in any
  CONSTRUCT output (encoded but never asked about), or vice versa.
  The protocol is then revised: either a question's scope is
  widened, or a new question is introduced and back-applied to all
  prior papers.
\end{itemize}

\paragraph{Limits and assumptions.}
Two assumptions worth flagging. First, the protocol assumes one
paper $\equiv$ one MLIP method $\times$ one material system. Papers
training multiple potentials or covering multiple systems
(e.g., Kumar et al.~2025 has separate C15 and C14 potentials)
instantiate the protocol per phase or per system, with the
appropriate identifiers (\texttt{kumar2025-c15},
\texttt{kumar2025-c14}). Second, papers occasionally report
ontology-relevant data the protocol does not capture (uncertainty
quantification, transferability claims, foundation-model fine-tuning
tags). When such data appears in three or more papers we treat that
as a signal to revise the protocol.

%% --------------------------------------------------------------------
%% Per-paper protocol-derived subsections begin here. Each is generated
%% from artifacts/kg/papers/<paper-id>.ttl + the 11 CONSTRUCT queries.
%% --------------------------------------------------------------------

\todo[inline]{Per-paper subsections to be added by the corpus-encoding
agent. Target: 19 papers covering the diversity targets in the
publication calendar.}

%% ====================================================================
